\section{Introduction}
\label{sec:intro}

\IEEEPARstart{C}{onvolutional} networks support a broad range of onboard
robotic perception workloads, from high-rate motion and geometry estimation to
detection and pose estimation~\cite{fastdepth,nanoflownet,bisenet,densefusion}.
Their outputs feed online navigation, planning, and manipulation, so inference
latency limits how often downstream components can incorporate new
observations~\cite{contactgraspnet}. Meeting this latency demand is difficult
on board-level processors with constrained compute and power, because dense
convolution recomputes every spatial position on every invocation. Prior
systems avoid part of this work by deriving active masks from event-camera
increments, temporal frame differences, or learned spatial
gates~\cite{evconv,asyncsparse,deltacnn,dgnet,dynconv}. Although these policies
arise from different sensing modalities and tasks, they expose the same
operator-level problem: update selected positions on dense 2D feature grids
without computing inactive outputs.

Dense GPU convolution obtains high throughput from tiled, coordinate-driven
execution. Each tile enumerates a contiguous block of output coordinates; those
coordinates determine the input reads and output writes, while the reduction
over channels and kernel offsets follows a regular datapath. Neighboring
coordinates consequently share predictable memory access and data reuse. This
regularity, not merely the arithmetic count, is what a masked operator must
preserve.

Existing masked-convolution systems sacrifice one part of this structure or the
other. \emph{Tile skipping}~\cite{sbnet,skipconv,evconv,deltacnn} retains the
dense pipeline and suppresses a tile only when all its outputs are inactive, so
a single active position triggers computation for the entire tile.
\emph{Gather-scatter}~\cite{sai,dgnet,dynconv} instead extracts active
positions exactly, materializes their receptive fields, and scatters the
results back to the dense grid, trading regular execution for extraction,
irregular access, and writeback.

We characterize this trade-off along two axes (\S\ref{sec:bg:gap}):
\emph{throughput} (operations per second, bounded by the dense pipeline) and
\emph{useful-compute ratio} (the fraction of executed operations that produce
needed outputs). Tile skipping preserves throughput but sacrifices useful-compute
ratio; gather-scatter attains a useful-compute ratio near one but sacrifices
throughput and pays for extraction and writeback. Neither delivers both, so
upstream sparsity does not translate into proportional latency reduction. In our
sparsest workload, the masks render 84\% of the convolution work they govern
unnecessary. State-of-the-art tile skipping~\cite{deltacnn} still spends half of
its executed operations on unneeded outputs (a useful-compute ratio of 0.49) and
cuts latency by at most 41\%. Gather-scatter drives the useful-compute ratio close
to one, yet sustains only 9\% of tile skipping's throughput and runs 3.1$\times$
slower than dense convolution (\S\ref{sec:eval:latency}).

Our key insight is that position-level selectivity does not require abandoning
the dense convolution pipeline. A dense tile already derives its reads and
writes from output coordinates, so replacing its contiguous coordinate block
with an active-output worklist changes which outputs are evaluated without
changing the per-position reduction, the dense tensor layout, or the weight
layout. Selectivity then becomes exact by construction, pinning the
useful-compute ratio near one with no input-patch materialization. A high
useful-compute ratio alone does not yield latency reduction, however, since
gather-scatter also attains one and still loses. What a worklist operator must
additionally control is the cost of constructing the worklist and the throughput
at which it is consumed.

Two challenges follow. First, \emph{worklist-construction overhead}: constructing the worklist requires scanning the mask, and incurs cost that scales with the layer's spatial grid, not with the
active count or the channel widths, while the compute the worklist enables shrinks with
both the active count and the channel widths. 
Construction's share of latency therefore grows as activity falls
and channels thin, and accumulates per operator when every layer rebuilds.
Second, \emph{end-to-end worklist efficiency}: the worklist's order and length both
affect its consumption throughput, and sparsity worsens each.
An unordered worklist scatters the input neighborhoods that adjacent work
touches, lowering locality; its length falls below the dense grid and shifts with
layer and input, so no fixed tiling keeps the GPU full. Both effects sharpen
as activity falls.

We present \sysname, a \emph{Worklist-driven Masked Convolution} operator that
answers both challenges by co-designing worklist construction and convolution.
Construction is made cheap and infrequent: a single mask-space pass fuses
output-mask propagation with per-tile compaction, moving only mask and
coordinate metadata, and compatible operators propagate and reuse valid
coordinate metadata instead of reconstructing it. Consumption is made
dense-like: construction emits contiguous per-tile segments whose order follows
the dense tile traversal, giving compute tile-local spatial structure without a
global sort, and compute organizes parallel work adaptively over the spatial and
reduction dimensions according to device, layer shape, and activity, while each
selected position retains the regular reduction over channels and kernel
offsets.

This paper makes three contributions:
\begin{itemize}
  \item We characterize masked-convolution latency along two axes,
        useful-compute ratio and throughput, and show that tile skipping and
        gather-scatter each secure one at the other's expense, explaining the
        measured gap between required work and realized latency.

  \item We design \sysname, which fuses position-level selectivity into tiled
        dense convolution to secure both axes at once: an exact worklist pins
        the useful-compute ratio near one, fused mask-space construction with
        cross-operator metadata reuse bounds construction cost, and tile-local
        segment ordering with adaptive work decomposition keeps throughput close
        to the dense pipeline.

  \item We evaluate \sysname on four perception models spanning event
        increments, temporal frame differences, and learned spatial gating,
        across two Jetson embedded platforms and discrete laptop and desktop
        GPUs. \sysname reduces full-model latency relative to the fastest baseline by
        22--39\% on the Jetson platforms and 34--46\% on the discrete GPUs, and
        cuts per-frame energy over dense convolution by 28--60\% on the Jetson
        platforms.
\end{itemize}
